\section{Discussion and conclusion}
For a fixed apparatus the important simplification is local in the material law. One common positive solve suffices for all interval changes at selected current and total-field rows. The calculation preserves each bound at the chosen weights. To change which media can be ruled out requires additional physical information. In particular, treating a known vacuum cell as an unknown current can hide a larger loss than the gain from imposing exact equalities on a few selected rows.

For spatial interpretation, the size of the specified inclusion matters. An array can reject most centres of a core with positive thickness while failing to reject even one isolated high-cell event with the rest of the medium free. A compatible centre carries a genuine full-wave witness. An undecided region may contain none. Uniform bounds on a finite-dimensional coefficient class in other inverse problems require their own assumptions on the class and on how the measurements are chosen \cite{alberti19,alberti21}. Here the array is prescribed, the physical class is explicit and the data can leave appreciable uncertainty.

The second result puts a further limit on active measurement. Saturation uses the signs available from the old data. The incident fields and the material form must also control all free receiving directions. At the native scale the two conditions can give different answers. On the compact examples, predictive bounds did separate material classes at fixed power, and uniform completion was already sufficient for the actual first extra exposure. A claim of improved reconstruction or of uniformly fewer microwave transmissions would require a different comparison. The separation margin instead provides an actionable decision: a particular physical beam and a specified readout tolerance suffice to resolve the stated alternatives.

The analysis assumes a common real scalar dielectric at a single frequency and a known linear receiving and illumination model. A complex material requires a dissipative inequality with the correct sign; a change of frequency also changes the Green operator. Unknown loaded-system error must enter the data ball before a result can be transferred to a physical specimen. The present calculations give constructive exclusion and possibility statements for finite full-wave models. They leave open how often larger material classes and broader acquisition uncertainties admit useful bounded predictions.

\section*{Data availability}
The source programs, synthetic arrays, measurement decisions and interval calculations accompany the research record. The two laboratory data collections are available from their original investigators \cite{fresnel,manitoba}. 
\section*{Acknowledgements}
OpenAI Codex assisted with mathematical development, implementation and drafting. Google Gemini assisted with independent examination of derivations and numerical evidence. The authors are responsible for the mathematics, numerical results and text.
